\section{Physical actions and the spectrum of reduced memory}
\label{nuclear:12}
\subsection{Genealogy of the dimensional sections}

The coordinates \(\pi,\varphi,e,\alpha\) are realizations of the correlated outputs constructed in Sections~\ref{sec:alpha} and~\ref{sec:electron}. This section uses \(\alpha=\alpha_A\), with the orientation and periodization of Definition~\ref{def:alpha-operativa}:

\begin{equation}
\alpha=\pi+e-\varphi-4-\kappa_{\rm per},\qquad
\kappa_{\rm per}=\frac{N_K}{1000^{12}-1},
\quad N_K=234543140729659824621058914794146601.
\label{nuc12:eq:2.1}
\end{equation}

Identity~\eqref{eq:alpha-identidad-completa} follows from the recurrence with compatible carry and its telescoping, proved in Proposition~\ref{prop:alpha-telescopia}. Period twelve corresponds to the register: a window of twelve blocks with a zero boundary condition and a compatible infinite section are different objects. Equality \eqref{nuc12:eq:2.1} uses the latter. The ordered register \(K\) preserves its incidences in addition to admitting the scalar reading \(\kappa_{\rm per}\).

The action reader, constructed in Section~\ref{iv:action:section}, determines

\begin{equation}
\begin{aligned}
H_5&=\tfrac12\sqrt{1000\alpha/\varphi}-\alpha+
\tfrac9{16}\alpha^2-\tfrac59\alpha^3+
\tfrac7{48}\alpha^4-\tfrac1{54}\alpha^5,\\
D_A&=\exp(-100\pi\alpha/9),\\
\mathcal C_\pi&=169\alpha^6+\frac{D_A\alpha^7}{1-D_A\alpha},\\
\eta_{\rm ret}&=H_5-\frac{90}{\pi}\mathcal C_\pi,\\
\hbar_{\rm pre}&=H_5\mathcal A_0,\qquad
\hbar_{\rm ret}=\eta_{\rm ret}\mathcal A_0,\qquad
R_{\rm act}=H_5/\eta_{\rm ret},
\end{aligned}
\label{nuc12:eq:2.2}
\end{equation}

Here \(\mathcal A_0=10^{-34}\mathcal U_S\). The decade comes from the determinantal reader of \(\varphi\alpha^{16}\), whose construction and bounds appear in Theorem~\ref{iv:action:decada}; \(H_5\) and \(\eta_{\rm ret}\) are the mantissas of the two positive sections of reduced action. In each section \(a\in\{\mathrm{pre},\mathrm{ret}\}\), we have \(h_a=2\pi\hbar_a\).

The construction of the central electron, developed in Section~\ref{sec:electron}, determines the coordinate

\begin{equation}
\begin{aligned}
\Delta_4&=\frac{\pi-e\log\pi}{270}\left(1+\frac{\pi}{729}\right),\\
B_e&=\frac{\sqrt3}{4}
\left[\exp(\varphi/\pi^2)-22\alpha^3\right](1+15\Delta_4),\\
\kappa_e&=80-54\alpha+6\Delta_4,\qquad
\varepsilon_e^\beta=B_eR_{\rm act}^{\kappa_e}.
\end{aligned}
\label{nuc12:eq:2.3}
\end{equation}

The prefactor \(\sqrt3/4\) comes from the norm of the commutator of the sum and product projectors on the fibres specified in \eqref{eq:el-prefactor}. The central route and its registers determine the return coefficients in \eqref{eq:el-kappa} and~\eqref{eq:el-beta}. The dimensional realization is \(E_e^\beta=\varepsilon_e^\beta\mathcal U_E\), \(m_e^\beta=E_e^\beta/c^2\), with the energy chart of \eqref{eq:el-dimensional}.

The dodecaphasic clock determines the energy basis of \eqref{eq:el-cuanto-reloj}:

\[
\omega_{108}=\frac{2\pi}{108t_0},\quad
E_{\rm clk}=\hbar_{\rm ret}\omega_{108},\quad
b_e=B_e\frac{\mathcal U_E}{E_{\rm clk}}.
\]

Substitution of \eqref{nuc12:eq:2.3}, preserving the same physical energy, proves

\begin{equation}
\boxed{\frac{m_e^\beta c^2}{\hbar_{\rm ret}}
=b_eR_{\rm act}^{\kappa_e}\omega_{108},\qquad
\frac{m_e^\beta c^2}{\hbar_{\rm pre}}
=b_eR_{\rm act}^{\kappa_e-1}\omega_{108}.}
\label{nuc12:eq:2.4}
\end{equation}

Division by \(h_a\), rather than \(\hbar_a\), converts angular frequency into cyclic frequency. The conversion \(\mathcal U_E/E_{\rm clk}\) remains explicit: a coordinate expressed in \(\mathrm{MeV}\) changes when its energy basis is replaced by the clock quantum.

\subsection{Normalization of the Maxwell--Dirac action}

The constitutive response in Section~\ref{iii:sec:constitutiva} produces the speed \(c\), impedance \(Z\), permittivity \(\varepsilon\), and permeability \(\mu\), all positive, with \(\varepsilon c=Z^{-1}\) and \(\mu c=Z\). The positive charge section in \eqref{iii:eq:charge} satisfies

\begin{equation}
e_a^2=\frac{2\alpha h_a}{Z}
=4\pi\alpha\varepsilon\hbar_a c .
\label{nuc12:eq:3.1}
\end{equation}

Here \(e_a>0\) denotes the elementary charge of section \(a\); \(e\) in \eqref{nuc12:eq:2.1}--\eqref{nuc12:eq:2.3} denotes the Euler coordinate. Fix a homogeneous chart with positive constants \(c,Z,\varepsilon,\mu,\hbar_a,e_a\), mass \(m=m_e^\beta>0\), and charge \(q\in\{-e_a,e_a\}\); then \(\sigma=q/e_a\in\{-1,1\}\). In the Minkowski realization of signature \((+,-,-,-)\), with \(x^0=ct\), the Maxwell--Dirac action is

\begin{equation}
\mathcal I=\int d^4x\left[
-\frac{F_{\mu\nu}F^{\mu\nu}}{4\mu c}
+\bar\psi\left(i\hbar_a\gamma^\mu
(\partial_\mu+iqA_\mu/\hbar_a)-mc\right)\psi
\right].
\label{nuc12:eq:3.2}
\end{equation}

The component \(A_0\) is the scalar potential divided by \(c\), and \(F=dA\). The spinor representation is fixed through the matrices \(g^I\) in \eqref{vii:eq:clifford-material-explicito}, which satisfy \(\{g^I,g^J\}=2\eta^{IJ}I_4\) for \(\eta=\operatorname{diag}(-1,1,1,1)\). Define \(\Gamma^I=-ig^I\), so that \(\{\Gamma^I,\Gamma^J\}=-2\eta^{IJ}I_4\), and, in this flat chart, \(\gamma^\mu=\delta^\mu_I\Gamma^I\), \(\bar\psi=\psi^\dagger\Gamma^0\). Signature, adjoint, and imaginary factor are thereby fixed jointly; the matrices \(\gamma^\mu\) are distinguished from the scalar Holst parameter used later. Consider smooth fields and compactly supported variations, or boundary conditions that cancel the boundary terms. The action realizes the already generated coefficients and the material representation of Section~\ref{vii:subsec:accion-espinorial-afin}.

Define

\[
\ell_a=\frac{\hbar_a}{mc},\quad X=x/\ell_a,\quad
\Psi=\ell_a^{3/2}\psi,\quad
\mathsf a_\mu=\frac{e_a\ell_a}{\hbar_a}A_\mu,\quad
\mathsf f=d\mathsf a.
\]

Then

\begin{equation}
\boxed{\frac{\mathcal I}{\hbar_a}=
\int d^4X\left[
-\frac{\mathsf f_{\mu\nu}\mathsf f^{\mu\nu}}{16\pi\alpha}
+\bar\Psi\bigl(i\gamma^\mu(\partial_\mu+i\sigma\mathsf a_\mu)-1\bigr)\Psi
\right].}
\label{nuc12:eq:3.3}
\end{equation}

\begin{proof}
The volume change contributes \(\ell_a^4\), each derivative contributes \(\ell_a^{-1}\), and the product of fermionic fields contributes \(\ell_a^{-3}\). The coefficient of the mass term is \(mc\ell_a/\hbar_a=1\). The electromagnetic term contributes \(\hbar_a/(4\mu c e_a^2)\). Since \(\mu c=Z\) and \(Ze_a^2=4\pi\alpha\hbar_a\), this coefficient is \(1/(16\pi\alpha)\). Both terms and their normalizations are thereby determined.
\end{proof}

Variation with respect to \(\mathsf a_\nu\), followed by substitution of \eqref{nuc12:eq:2.1}, gives

\begin{equation}
\partial_\mu\mathsf f^{\mu\nu}
=4\pi\left(\pi+e-\varphi-4-\kappa_{\rm per}\right)
\sigma\bar\Psi\gamma^\nu\Psi .
\label{nuc12:eq:3.4}
\end{equation}

The cancellation composes the action sections, electronic construction, and constitutive response; the regional and dodecaphasic reader determines the remaining coupling. Mass persists in the length \(\ell_a\), and action in the dimensional conversion. Compton normalization is a subsequent realization of these outputs. In the Coulomb realization of the same electron, the Bohr length is \(\ell_a/\alpha\): substituting \(e_a^2=4\pi\alpha\varepsilon\hbar_a c\) into \(4\pi\varepsilon\hbar_a^2/(m e_a^2)\) gives exactly this expression.

The deduction uses fixed homogeneous coefficients. For spacetime-dependent coefficients, the product rule introduces derivatives of the scales and coupling into the transformed equations. Variation of an action with variable parameters must preserve these terms.

\subsection{Areal scale of gravitational action and entropy}

The circular realization of the \(108\)-step cycle, constructed in Section~\ref{v:ant:sec:gravity}, transports the length \(L=L_*\alpha^{16}5759/23040\) from~\eqref{xii:radial:length}. With \(t_0=\pi L/(54c)\), it determines

\begin{equation}
t_P=\frac{54}{\pi}t_0,\qquad
G_a=\frac{c^5t_P^2}{\hbar_a},\qquad
\ell_P^2=\frac{\hbar_aG_a}{c^3}=c^2t_P^2.
\label{nuc12:eq:4.1}
\end{equation}

The radius-matching condition of this realization is \(G_a m_{108,a}/c^2=R_{108}=L\), with \(m_{108,a}=\hbar_a\omega_{108}/c^2\). It is the restriction of~\eqref{xii:radial:G} to the unit mode; the chronological form retains the factor \((54/\pi)^2\). The reader uses the reduced radius \(G_a m/c^2\). Reciprocity \(G_{\rm pre}/G_{\rm ret}=R_{\rm act}^{-1}\) preserves \(\ell_P^2=L^2\) between the two sections. The reference \(L_*\), edge base \(u_L\), and advance \(ct_0\) are related by~\eqref{xii:radial:clock}.

Section~\ref{vii:sec:realizacion-cartan} constructs the cotetrad and connection from the transport already realized. In the action of Section~\ref{vii:sec:corriente-respuesta-cartan}, fix a nondegenerate cotetrad \(e^I\), the orientation, signature \(\eta=\operatorname{diag}(-1,1,1,1)\), and a metric connection \(\omega\), with curvature \(F^{IJ}=d\omega^{IJ}+\omega^I{}_{K}\wedge\omega^{KJ}\). Take \(\gamma\in\mathbb R\setminus\{0\}\) and hold \(\gamma,\Lambda,c,G_a,\hbar_a\) constant during variation. Put \(\Sigma^{IJ}=e^I\wedge e^J\) and \(\mathsf P_\gamma=\star+\gamma^{-1}I\); internal Lorentzian duality satisfies \(\star^2=-I\) on bivectors. Variations have compact support or boundary conditions that cancel the surface term. Substitution of \eqref{nuc12:eq:4.1} into \eqref{vii:eq:accion-cartan-holst} gives

\begin{equation}
\boxed{\frac{\mathcal I_{\rm grav}}{\hbar_a}
=\frac1{16\pi\ell_P^2}
\left[\int\Sigma_{IJ}\wedge(\mathsf P_\gamma F)^{IJ}
-2\Lambda\int\operatorname{vol}_e\right].}
\label{nuc12:eq:4.2}
\end{equation}

With the cotetrad and material fields fixed during variation of the connection, let \(\sigma_{IJ}\) be the material current defined by \(\delta_\omega\mathcal I_m=-\tfrac12\int\delta\omega^{IJ}\wedge\sigma_{IJ}\), according to \eqref{vii:eq:corriente-variacional}. The same substitution in Theorem~\ref{vii:thm:ecuacion-cartan-holst} produces

\begin{equation}
D_\omega(\mathsf P_\gamma\Sigma)
=8\pi\ell_P^2\,\frac{\sigma}{\hbar_a}.
\label{nuc12:eq:4.3}
\end{equation}

\begin{proof}
In section \(a\), the original coefficient is \(1/(2\kappa_a c)\), with \(\kappa_a=8\pi G_a/c^4\). Division by \(\hbar_a\) converts it into \(1/(16\pi\ell_P^2)\); the cosmological term retains its factor two. In the connection equation, \(\kappa_a c\hbar_a=8\pi\ell_P^2\) holds. Both identities use the same dimensional section.
\end{proof}

The Barbero functional of Section~\ref{v:ant:sec:barbero} retains its complete expression. Its channels are \(q_\pm=e^{-x\mp y}\), with \(x=\pi A/180\), \(y=\pi C^*/180\), where \(A\) and \(C^*\) are the angular coordinates expressed in degrees. The domain \(x>|y|\) ensures \(0<q_\pm<1\):

\begin{equation}
\gamma=\frac{180}{\pi}xy+
12[S_{90}(q_-)-S_{90}(q_+)]
-[S_{120}(q_-)-S_{120}(q_+)],\qquad
S_m(q)=\frac{q^m}{1-q^{3m}}.
\label{nuc12:eq:4.4}
\end{equation}

Hexadic and octadic incidence fixes the multiplicities \(90\) and \(120\), and the oriented chain fixes the coefficients \(12\) and \(-1\), as proved by \eqref{v:ant:eq:gamma}. The areal construction of Section~\ref{v:ant:sec:area} additionally gives

\begin{equation}
a_0=\frac{1+2\cos(\pi/18)}3\,\frac y6,\qquad
\widehat{\mathcal A}=\gamma a_0\ell_P^2(N_{90}+N_{120}).
\label{nuc12:eq:4.5}
\end{equation}

The inversion \(y\mapsto-y\) changes the signs of \(\gamma\) and \(a_0\); their product is even. The operator \(\mathsf P_\gamma\), by contrast, retains its orientation dependence. Invariance of the areal product is thus distinguished from invariance of each operator coefficient. The areal spectrum is realized in the positive section \(\gamma a_0>0\), with the occupation operators \(N_{90},N_{120}\geq0\) constructed in Section~\ref{v:ant:sec:area}.

For the spherical cosmological horizon in \eqref{vii:eq:horizonte-cosmologico}, of radius \(R>0\) and area \(\mathcal A_H=4\pi R^2\), fix a section \(a\) and abbreviate \(\hbar=\hbar_a\), \(G=G_a\) during this calculation and the following bounce calculation. Then

\begin{equation}
\frac{S_H}{k_B}=\frac{\mathcal A_H}{4\ell_P^2},\qquad
T_{\rm cos}=\frac{\hbar c}{2\pi k_BR},\qquad
E_\Lambda=\frac{c^4R}{2G},\qquad T_{\rm cos}S_H=E_\Lambda .
\label{nuc12:eq:4.6}
\end{equation}

The areal scale in \eqref{nuc12:eq:4.2}, \eqref{nuc12:eq:4.3}, and \eqref{nuc12:eq:4.6} is the same. The temperature corresponds to the spherical cosmological horizon; the normalization of the Schwarzschild horizon is a different realization. Varying \(R\) while holding \(c,G,\hbar,k_B\) fixed, differentiation of \eqref{nuc12:eq:4.6} retains both terms: \(T_{\rm cos}\,dS_H=2\,dE_\Lambda\) and \(S_H\,dT_{\rm cos}=-\,dE_\Lambda\), as in \eqref{vii:eq:diferencial-horizonte}. This gives the differential of an identity between states in that collection, not a replacement for the dynamical balance with other variable parameters.

\subsubsection{An explicit dynamical application}

Theorem~\ref{vii:thm:amplitud-material-explicita} determines the following amplitude for a spinor state with conserved second moment \(q_\varrho>0\). Conservation is expressed through \eqref{vii:eq:conservacion-momento-cuartico}; Corollary~\ref{vii:cor:dominio-positivo-conservado} provides a concrete realization in the three-occupation sector:

\begin{equation}
s_0=\frac{3\pi G\hbar^2}{2c^2V_c^2}
\frac{\gamma^2}{1+\gamma^2}q_\varrho .
\label{nuc12:eq:4.7}
\end{equation}

In the spatially flat realization with pressureless matter and \(\Lambda=0\), the constant \(\rho_0>0\) is an energy density, \(V_c>0\) is the reference comoving volume, and \(E_0=\rho_0V_c>0\). The solution of Theorem~\ref{vii:thm:rebote-polvo} is

\[
a_b^3=s_0/\rho_0,\qquad
\tau^2=\frac{s_0c^2}{6\pi G\rho_0^2},\qquad
a(t)=a_b\left(1+\frac{(t-t_b)^2}{\tau^2}\right)^{1/3}.
\]

Restoring \(c\) in \eqref{vii:eq:dinamica-homogenea} uses the coefficient \(8\pi G/(3c^2)\) when the density is an energy density. In particular, \((\dot a/a)^2=(8\pi G/(3c^2))(\rho_0a^{-3}-s_0a^{-6})\). Substitution of \eqref{nuc12:eq:4.7} into the solution parameters proves

\begin{equation}
\boxed{\left(\frac{\tau E_0}{\hbar}\right)^2
=\frac{q_\varrho}{4}\frac{\gamma^2}{1+\gamma^2},\qquad
\mathcal V_b=\frac{3\pi}{2}q_\varrho
\frac{\gamma^2}{1+\gamma^2}\ell_P^2\frac{\hbar c}{E_0}.}
\label{nuc12:eq:4.8}
\end{equation}

Here \(\mathcal V_b=V_c a_b^3\) is the physical volume at the bounce. The quantities \(G\) and \(V_c\) cancel in the first equality; \(E_0\) remains as material energy. In the conserved sector \(\Lambda^3\mathbb C^4\), Proposition~\ref{vii:prop:operador-espin-car} proves \(\widehat{\mathscr Q}_{\rm sp}=4I\), so \(q_\varrho=4\) for every normalized state supported in that sector, and \(\tau E_0/\hbar=|\gamma|/\sqrt{1+\gamma^2}\). This yields a dynamical scale for the specified spinor state and cosmological realization.

\subsection{Prolongation of memory transport: quantum realization}

The quantum realization uses the reversible coupling of Section~\ref{nuclear:10}. First consider a real canonical plane with alternating matrix \(\Omega=\left(\begin{smallmatrix}0&1\\-1&0\end{smallmatrix}\right)\). The nonadic partition determines

\begin{equation}
Q=\frac13\begin{pmatrix}\sqrt8 I&-I\\I&\sqrt8 I\end{pmatrix},\qquad
U_H=Q^{\mathsf T}\operatorname{diag}(I,H)Q
=\frac19\begin{pmatrix}8I+H&\sqrt8(H-I)\\
\sqrt8(H-I)&I+8H\end{pmatrix}.
\label{nuc12:eq:5.1}
\end{equation}

The operator \(H\in\operatorname{Sp}(2,\mathbb R)\) transports this plane. The matrix \(Q\) preserves the Euclidean inner product and the form \(\Omega\oplus\Omega\); hence \(U_H\) is symplectic. The first component defines the observed system and the second the memory register; the complete transformation acts on both.

The ordered loop constructed in Section~\ref{nuclear:11} is

\begin{equation}
C=\begin{pmatrix}0&-1\\1&-1\end{pmatrix},\quad
P=\begin{pmatrix}1&1\\0&1\end{pmatrix},\quad
H_n=C^{-1}P^{-n}CP^n
=\begin{pmatrix}1+n&n^2\\n&n^2-n+1\end{pmatrix}.
\label{nuc12:eq:5.2}
\end{equation}

For every \(n\in\mathbb Z\), \(\det H_n=1\) holds. In particular, for \(n=1\),

\begin{equation}
H_1=\begin{pmatrix}2&1\\1&1\end{pmatrix}
=F_1^2,\quad F_1=\begin{pmatrix}1&1\\1&0\end{pmatrix}
=P F_{\rm av}P^{-1}.
\label{nuc12:eq:5.3}
\end{equation}

Its spectrum is \(\{\varphi^2,\varphi^{-2}\}\), by subsequent recognition of the already generated golden mode. The integer \(n\) specifies the oriented number of iterations in the chosen loop. The transport collection is determined algebraically; its realization as a physical protocol additionally preserves the corresponding preparation and evolution data.

\subsubsection{Initial state, transport, and reading}

Proposition~\ref{prop:ii-covarianza-gaussiana} constructs the Gaussian realization of the action ellipse. In the positive section \(\hbar_a>0\), with \(A>0\) and \(|C^*/A|<1\), its covariance is

\begin{equation}
V_\eta=\frac{\hbar_a}{2}M_\eta M_\eta^{\mathsf T},\qquad
M_\eta=\operatorname{diag}(e^{\eta/2},e^{-\eta/2}),\qquad
\eta=\operatorname{artanh}(C^*/A).
\label{nuc12:eq:5.4}
\end{equation}

In the Schrödinger representation on \(L^2(\mathbb R,dQ)\), with \(\widehat Q\) multiplication by \(Q\) and \(\widehat P=-i\hbar_a\,d/dQ\), the relation \([\widehat Q,\widehat P]=i\hbar_a I\) holds on Schwartz space. The normalized pure state is

\[
\psi_\eta(Q)=(\pi\hbar_a e^\eta)^{-1/4}
\exp[-Q^2/(2\hbar_a e^\eta)].
\]

The preparation is the product \(\psi_\eta\otimes\psi_\eta\). The transport in \eqref{nuc12:eq:5.1} is expressed in the same covariance chart by

\begin{equation}
\widetilde U=(M_\eta\oplus M_\eta)U_H
(M_\eta^{-1}\oplus M_\eta^{-1}).
\label{nuc12:eq:5.5}
\end{equation}

Lemma~\ref{app:lem:implementacion-simplectica} provides a unitary implementation of this symplectic matrix in two modes. Implementations differing by a global phase produce the same conjugation of the state. Covariance transforms by \(V\mapsto\widetilde U V\widetilde U^{\mathsf T}\); the joint state remains pure, and its partial trace over the second mode determines the visible reading, according to the construction in Section~\ref{app:gaussianas-espectro}.

\begin{theorem}[Reduced mixedness from relative deformation]

The covariance of the visible reading is

\begin{equation}
\boxed{V_{\rm vis}=
\frac{\hbar_a}{2}M_\eta
\frac{8I+HH^{\mathsf T}}9
M_\eta^{\mathsf T}.}
\label{nuc12:eq:5.6}
\end{equation}

Its normalized symplectic eigenvalue is

\begin{equation}
\boxed{\nu(H):=\frac{2\sqrt{\det V_{\rm vis}}}{\hbar_a}
=\sqrt{1+\frac8{81}\left[\operatorname{tr}(HH^{\mathsf T})-2\right]}.}
\label{nuc12:eq:5.7}
\end{equation}
\end{theorem}

\begin{proof}
Let \(T=(8I+H)/9\) and \(D=\sqrt8(H-I)/9\). The two initial modes are independent and have the same covariance. In the balanced chart of the initial state, the normalized visible covariance is \(TT^{\mathsf T}+DD^{\mathsf T}\). Expansion gives

\[
\begin{aligned}
TT^{\mathsf T}+DD^{\mathsf T}
&=\frac{64I+8(H+H^{\mathsf T})+HH^{\mathsf T}
+8[HH^{\mathsf T}-H-H^{\mathsf T}+I]}{81}\\
&=\frac{8I+HH^{\mathsf T}}9.
\end{aligned}
\]

This proves \eqref{nuc12:eq:5.6}. Since \(\det M_\eta=1\) and \(\det(HH^{\mathsf T})=1\), \(\det(8I+HH^{\mathsf T})=65+8\operatorname{tr}(HH^{\mathsf T})\). Division by \(81\) proves \eqref{nuc12:eq:5.7}. The two positive eigenvalues of \(HH^{\mathsf T}\) are reciprocal; their sum is at least two. Consequently, \(\nu\geq1\), with equality exactly when \(HH^{\mathsf T}=I\).
\end{proof}

The scale \(\hbar_a\) and deformation \(\eta\) are retained until the determinant is evaluated. Their cancellation is covariant because state and operator are transported jointly. The expression \(HH^{\mathsf T}\) uses the Euclidean inner product of the initial state's balanced chart; a chart change also transports this metric inner product.

An orthogonal symplectic transport \(H\) may have \(D\neq0\) while preserving the isotropic Gaussian product, with \(\nu=1\). The operator complement and state entanglement therefore have distinct criteria. In this preparation, non-isometric relative deformation characterizes exactly the case \(\nu>1\).

\subsubsection{Exact evaluation of the golden loop}

For \eqref{nuc12:eq:5.2},

\begin{equation}
\operatorname{tr}(H_nH_n^{\mathsf T})-2
=n^2(2n^2-2n+5),\qquad
\nu_n^2=1+\frac8{81}n^2(2n^2-2n+5).
\label{nuc12:eq:5.8}
\end{equation}

The identity follows by summing the squares of the four entries. For \(n=1\), the sum is seven and

\begin{equation}
\boxed{\nu_1=\frac{11}{9},\quad
\operatorname{Tr}\rho_{\rm vis}^2=\frac9{11},\quad
\bar n_{\rm W}=\frac19,\quad
\frac{S_{\rm vis}}{k_B}=\frac{10}{9}\log10-\log9.}
\label{nuc12:eq:5.9}
\end{equation}

Theorem~\ref{app:thm:espectro-gaussiano} proves that a centred one-mode Gaussian state with symplectic eigenvalue \(\nu\) is unitarily equivalent to the geometric density operator in \eqref{app:eq:estado-geometrico}, whose mean occupation is \(\bar n_{\rm W}=(\nu-1)/2\). For \(\nu=11/9\), its eigenvalues are

\begin{equation}
\lambda_j=\frac9{10}\,10^{-j},\qquad j=0,1,2,\ldots.
\label{nuc12:eq:5.10}
\end{equation}

The series sums to one. The sum of its squared terms is \((9/10)^2/(1-10^{-2})=9/11\). Finally,

\[
-\sum_j\lambda_j\log\lambda_j
=-\log(9/10)+\log10\sum_jj\lambda_j
=\frac{10}{9}\log10-\log9.
\]

The Gaussian reduction and its spectrum are proved in Section~\ref{app:gaussianas-espectro}; the two preceding series evaluate exactly the purity and entropy of this realization. Since the joint state is pure, \(S_{\rm vis}/k_B\) is also its dimensionless entanglement entropy. The global transformation is unitary, whereas restriction to the first mode has positive entropy.

The result corresponds to the specified preparation and transport. Identifying it with a vacuum measurement requires experimental realization of that protocol. The entropy of a reduction and transferred heat remain distinct quantities.

\subsubsection{Composition with the thermal transducer}

In the Williamson normal representation, choose a frequency \(\omega>0\) and define the energy reader \(\widehat H_{\rm W}=\hbar_a\omega(\widehat N+\tfrac12)\). The frequency is part of this energy realization; covariance determines the populations. With this choice, \eqref{nuc12:eq:5.10} is a Gibbs state whose population ratio satisfies

\[
e^{-\hbar_a\omega/(k_BT_{\rm W})}=\frac1{10},\qquad
T_{\rm W}=\frac{\hbar_a\omega}{k_B\log10}.
\]

If the clock frequency of the same section is taken, \(\omega=1/t_P=2\pi/T_{108}\), the thermal-transducer identity \eqref{v:ant:eq:ii-kb-aut-par-termico}, \(\hbar_a/t_P=k_B\Theta_{{\rm clk},a}\log3\), gives

\begin{equation}
\boxed{\frac{T_{\rm W}}{\Theta_{{\rm clk},a}}=\frac{\log3}{\log10}.}
\label{nuc12:eq:5.11}
\end{equation}

The action section and thermal transducer cancel in the ratio. The index \(a\) of \(\Theta_{{\rm clk},a}\) preserves compatibility with the section \(\hbar_a\), as in \eqref{v:ant:eq:ii-kb-aut-orientaciones}. The temperature \(T_{\rm W}\) is the Gibbs-equivalent temperature for the declared Hamiltonian; the isolated unitary evolution retains its reversible character. This spectral identification determines a thermal ratio and does not postulate thermalization dynamics.

\subsection{Extension to arbitrary finite canonical dimension}

Let \(H\in\operatorname{Sp}(2d,\mathbb R)\), with finite \(d\geq1\). Prepare two pure, centred, identical Gaussian states of \(d\) modes each, with covariance \(V_0=(\hbar_a/2)MM^{\mathsf T}\), where \(M\in\operatorname{Sp}(2d,\mathbb R)\). Section~\ref{app:gaussianas-espectro} specifies this preparation and the joint chart transport by \(M\oplus M\). In the balanced chart, the same coupling \(U_H\) produces \(W=(8I+HH^{\mathsf T})/9\). The product \(A=HH^{\mathsf T}\) is positive, symmetric, and symplectic.

There exists an orthonormal symplectic basis in which

\[
A=\operatorname{diag}(e^{2r_1},e^{-2r_1},\ldots,
e^{2r_d},e^{-2r_d}),\qquad r_j\ge0.
\]

The identity \(A\Omega A=\Omega\), with \(\Omega=\bigoplus_{j=1}^d\left(\begin{smallmatrix}0&1\\-1&0\end{smallmatrix}\right)\), implies that \(\Omega\) takes the eigenspace of \(\lambda\) to that of \(\lambda^{-1}\). Paired orthonormal bases are chosen in these spaces. The eigenspace of \(\lambda=1\) is \(\Omega\)-invariant and admits orthonormal canonical pairs. Their union constructs the required basis.

The matrix \(W\) is diagonal in that basis. The products of each pair determine its symplectic eigenvalues:

\begin{equation}
\boxed{\nu_j^2=
\frac{(8+e^{2r_j})(8+e^{-2r_j})}{81}
=1+\frac{32}{81}\sinh^2r_j.}
\label{nuc12:eq:6.1}
\end{equation}

Theorem~\ref{app:thm:espectro-gaussiano} composes partial-trace reduction with these symplectic eigenvalues:

\begin{equation}
\boxed{\operatorname{Tr}\rho_{\rm vis}^2=\prod_{j=1}^d\nu_j^{-1},
\qquad
\frac{S_{\rm vis}}{k_B}=\sum_{j=1}^d
\left[\frac{\nu_j+1}{2}\log\frac{\nu_j+1}{2}
-\frac{\nu_j-1}{2}\log\frac{\nu_j-1}{2}\right].}
\label{nuc12:eq:6.2}
\end{equation}

Adopt \(0\log0=0\) by continuity. The proof of Theorem~\ref{app:thm:espectro-gaussiano} unitarily transforms the state into a product of \(d\) geometric density operators: the purity of the product is the product of the purities, and its entropy is the sum of the entropies. The result holds for every finite \(d\) and every transport \(H\) in the domain, retaining all occupations of Fock space.

For the loop \eqref{nuc12:eq:5.3}, \(r=2\log\varphi\) and \(\sinh^2r=5/4\), reproducing \eqref{nuc12:eq:5.9}. Thus the transport spectrum and nonadic partition \(8+1\) determine the spectrum of the reduced reading; the action sections fix its dimensional realization.

The dimension of an exceptional realization and its transport come from the corresponding constructor. Equation \eqref{nuc12:eq:6.1} determines the spectral law after realizing its canonical modes and the specified Gaussian preparation.

\subsubsection{Operator identity between memory and complex structure}

In the balanced canonical chart, let \(J=\Omega\), with \(J^2=-I\) and \(J^{\mathsf T}=-J\). This matrix realizes the canonical complex structure of phase space, and \(D=\sqrt8(H-I)/9\) is the memory block. For the same preparation and chart transport as in Section~\ref{app:gaussianas-espectro}, we have

\begin{equation}
\boxed{-(JW)^2=I+[D,J][D,J]^{\mathsf T},\qquad
W=\frac{2}{\hbar_a}M^{-1}V_{\rm vis}M^{-\mathsf T}.}
\label{nuc12:eq:6.3}
\end{equation}

\begin{proof}
The symplectic identities \(HJH^{\mathsf T}=J\) and \(H^{\mathsf T}JH=J\), together with \(A=HH^{\mathsf T}\), give \(JAJ^{\mathsf T}=A^{-1}\). If \(C=HJ-JH\), multiplication of the four summands produces

\[
CC^{\mathsf T}
=A+JAJ^{\mathsf T}
-HJH^{\mathsf T}J^{\mathsf T}
-JHJ^{\mathsf T}H^{\mathsf T}
=A+A^{-1}-2I.
\]

On the other hand,

\[
-(JW)^2
=\frac{(8I+A^{-1})(8I+A)}{81}
=I+\frac8{81}(A+A^{-1}-2I).
\]

Since \([D,J]=\sqrt8\,C/9\), we obtain \eqref{nuc12:eq:6.3}. The identity holds in every finite canonical dimension and for every \(H\) in the domain, including transports that mix the modes.
\end{proof}

The eigenvalues of this operator are the \(\nu_j^2\), each with multiplicity two. Formula~\eqref{app:eq:pureza-general} then expresses purity directly in terms of the memory block:

\begin{equation}
\boxed{\operatorname{Tr}\rho_{\rm vis}^2
=\det\!\left(I+[D,J][D,J]^{\mathsf T}\right)^{-1/4}.}
\label{nuc12:eq:6.4}
\end{equation}

In one mode,

\begin{equation}
\nu^2=1+\frac12\|[D,J]\|_F^2.
\label{nuc12:eq:6.5}
\end{equation}

In this protocol, \(D\) commutes with \(J\) exactly when the reduced reading preserves purity. For the loop \(H_1\), \(\|[D,J]\|_F^2=80/81\), and \eqref{nuc12:eq:6.5} gives \(\nu^2=121/81\), yielding the purity \(9/11\).

The commutator uses the complex structure compatible with the initial covariance and its metric inner product. A canonical change jointly transports all three data; the transpose in \eqref{nuc12:eq:6.3} refers to the balanced chart. The prior construction of the imaginary element in the formal kernel precedes this representation, which specifies its canonical action.

The joint evolution preserves purity through its unitary implementation; the commutator determines the loss of purity of the partial reading. The information of the joint state remains in the two factors and their correlations. The formula determines the purity of the reduction, which is a nonlinear functional of its density operator.

The real decomposition relative to \(J\) distinguishes its complex-linear and complex-antilinear components:

\[
D_{\rm lin}=\frac{D-JDJ}{2},\qquad
D_{\rm anti}=\frac{D+JDJ}{2}.
\]

The first component commutes with \(J\) and the second anticommutes with \(J\). The identities \([D,J]=2D_{\rm anti}J\) and \(JJ^{\mathsf T}=I\) imply

\begin{equation}
\boxed{-(JW)^2=I+4D_{\rm anti}D_{\rm anti}^{\mathsf T},\qquad
\operatorname{Tr}\rho_{\rm vis}^2
=\det(I+4D_{\rm anti}D_{\rm anti}^{\mathsf T})^{-1/4}.}
\label{nuc12:eq:6.6}
\end{equation}

The complex-antilinear component of the memory block, defined relative to the complex structure of the initial state, determines exactly the purity and mixedness spectrum of this protocol. The complex-linear component can transport memory while preserving purity. The distinction concerns linearity relative to \(J\); combinatorial complexity and the arithmetic classification of constants are properties of other objects.
